% ---
% title: All four suits
% tags: equally-likely
% ---

\begin{problem}
	Five cards are dealt from a well-shuffled deck of 52 cards (without
	replacement). The deck has four suits, with 13 cards in each suit.
	(a) What is the probability that all four suits appear among the five
	cards?
	(b) A classmate argues as follows: ``Choose one card of each suit
	($13^4$ ways), then any one of the remaining 48 cards. So the probability
	is $13^4\cdot48/\binom{52}{5}$.'' Is this argument correct? If not, what
	goes wrong?
\end{problem}

\begin{solution}
	\idea All hands of five cards are equally likely, so the probability is
	the number of good hands divided by the number of all hands. To count the good hands, first look at what such a hand must look like.
	\step{(a)} The sample space $\Om$ is the set of all 5-card hands. The order
	of the cards does not matter, so
	\[ |\Om|=\binom{52}{5}=2\,598\,960. \]
	Every hand is equally likely, because the deck is well shuffled.

	Now the good hands. There are 5 cards but only 4 suits. If all four suits
	appear, each suit has at least one card, and one card is left over. So
	exactly one suit appears twice and the other three suits appear once each.
	We build such a hand in three steps:
	\begin{itemize}
		\item choose the suit that appears twice: $4$ ways;
		\item choose its 2 cards out of its 13 cards: $\binom{13}{2}=78$ ways;
		\item choose one card from each of the other three suits: $13\cdot13\cdot13=13^3=2197$ ways.
	\end{itemize}
	Each good hand is built in exactly one way (the doubled suit is determined
	by the hand). By the multiplication rule the number of good hands is
	\[ 4\cdot\binom{13}{2}\cdot13^3=4\cdot78\cdot2197=685\,464. \]
	Hence
	\[ P(\text{all four suits})=\frac{4\cdot\binom{13}{2}\cdot13^3}{\binom{52}{5}}
		=\frac{685\,464}{2\,598\,960}=\frac{2197}{8330}\approx0.2637. \]
	\step{(b)} The argument is not correct: it counts every good hand twice.
	Take, for example, the hand
	\[ \heartsuit\mathrm{A},\ \heartsuit\mathrm{K},\ \spadesuit 5,\ \diamondsuit 9,\ \clubsuit 2. \]
	The classmate's procedure produces it in two different ways: once with
	$\heartsuit\mathrm{A}$ as ``the heart chosen first'' and $\heartsuit\mathrm{K}$ as ``the
	extra card'', and once with the roles of these two cards swapped. The same
	happens for every good hand: either card of the doubled suit can play the
	role of ``the chosen one''. So
	\[ 13^4\cdot48=1\,370\,928=2\cdot685\,464, \]
	exactly twice the correct count, and the classmate's probability
	$1\,370\,928/2\,598\,960\approx0.5275$ is twice the right value.
	(Dividing by 2 would repair the argument.)
	\begin{answer}
		(a) $\dfrac{2197}{8330}\approx0.2637$ \quad (b) No: each good hand is
		counted twice, so the result is twice the correct probability.
	\end{answer}
	\begin{remark}
		\textbf{Common mistake:} a ``choose a special element first, then the
		rest'' count often produces the same set in several ways. A quick test:
		take one concrete good outcome and ask in how many ways the procedure
		can produce it. If the answer is not 1, the count is wrong (here it is
		2 for every hand, so divide by 2).
	\end{remark}
\end{solution}
